% !TEX TS-program = lualatex
Диаграмма \refimage{model-evaluator-usecase} описывает взаимодействие пользователя с модулем для тестирования натренированных интеллектуальных агентов --- инструментом для оценки и сравнения ранее обученных моделей. Данный модуль предоставляет удобный интерфейс для загрузки моделей, проведения тестовых игр и вывода аналитических результатов. Участниками диаграммы являются пользователь и две обученные модели (Model 1 и Model 2), которые рассматриваются как внешние акторы, взаимодействующие с системой при загрузке.

\image[width=0.9\textwidth]{Схема использования модуля для тестирования натренированных интеллектуальных агентов}{model-evaluator-usecase}{model-evaluator-usecase}

Процесс начинается с запуска пользователем оценщика моделей. После этого пользователь указывает параметры оценки, среди которых пути к моделям (\texttt{-{}-}first, \texttt{-{}-}second), количество тестовых эпизодов (\texttt{-{}-}episodes), а также, при необходимости, параметры архитектуры скрытых слоёв (\texttt{-{}-}h1, \texttt{-{}-}h2). Параметр \texttt{-{}-}second является необязательным; его наличие определяет, будет ли сравнение происходить между двумя моделями или между одной моделью и случайным агентом.

После указания параметров производится загрузка первой модели. Если в параметрах указан \texttt{-{}-}second, дополнительно загружается вторая модель. Каждая из моделей участвует в соответствующем варианте использования как внешний актор, предоставляющий данные при инициализации.

Следующим шагом является создание игрового окружения, в котором будет происходить тестирование. После подготовки среды начинается основное тестирование --- проведение серии игр. В зависимости от наличия одной или двух моделей, система либо сравнивает их между собой, либо оценивает одну модель против случайного агента. В обоих случаях используется метод compare\_with, реализованный в базовом классе BaseModel.

По завершении серии игр выполняется вывод результатов. На этом этапе система предоставляет пользователю статистику, включающую число побед и общее количество сыгранных игр.
